\section{Lesson 3, 30/09}
\subsection{Equilibrium in external field}
The TD system is often subject to external forces, like gravity: how can equilibrium
be described in this case?

Consider a TD system and subdivide it into smaller systems s.t they are still
TD systems with a large number of particles: the coordinates of the portion is
$\mathbf{x}$.

Since $\mu$ is the energy needed to add particles to the system, the control
variables used are $T,p,N$: from the definition of $G$
$$\mu = \frac{G}{N}$$
and from the fundamental equations
$$\pdv[]{\mu}{p} = \frac{1}{N}\pdv[]{G}{p}:= \frac{V(T,p,N)}{N} = v(T,p)$$
where $v$ is the volume per particle.

The chemical potential consists of two parts
$$\mu(T,p;\mathbf{x}) = \mu_0(T,p)+u(\mathbf{x})$$
where $\mu_0$ is $\mu$ with the external field turned off and $u$ is related to the
field.

At the equilibrium, one expects that $T$ remains constant, $p$ depends on
$\mathbf{x}$ (like the Stevino law) and $\mu$ remains constant and homogeneous.

These conditions translate to
$$\pdv{\mu}{x^i} = 0 = \pdv{\mu}{p}\cdot\pdv{p}{x^i}+\pdv{u}{x^i}$$
so
$$v(T,p)\pdv{p}{x^i}+\pdv{u}{x^i}.$$
\textbf{Example: gravity} In this case $u(\mathbf{x})=mgz$, so
$$\dv[]{p}{z}v(T,p)=-mg$$
so $v$ does not depend on $T$, adding the incompressibility of the fluid imposes
constant $v$, so
$$\dv[]{p}{z} = -\frac{m}{v}gz=-\rho gz \rightarrow p(z) = p_0-mgz$$
in the case of an ideal gas, using $pv=k_bT$ yields (prove it as an exercise)
$$p(T,z)=p_0\exp\left[-\frac{mg}{kT}z\right]$$
\subsection{Second order equilibrium}
It was found that, at first order
\begin{align*}
&T,p,\mu \text{ are homogeneous}\\
&\Delta S = 0
\end{align*}
To discuss the second order fluctuations, it is needed to re-express everything in
terms of TD potentials: consider the control variables
$$\left\{\underbrace{X_1,...X_r}_\text{extensive};
\underbrace{I_{r+1},...I_n}_\text{intensive}\right\}$$
and the potential
$$\Phi(X_i;I_j):$$
its differential is
$$\dd\Phi = \sum_{i=1}^{r}I_i\dd X_i - \sum_{j=r+1}^{n}X_j\dd I_j$$
\textbf{Note} This generalizes the problem to have $r$ extensive control variables
and $r-n$ (?) intensive control variables.

Consider a TD system subdivided in two parts and a fluctuation in the extensive
variables:
$$\delta X_i = \delta X_{i,1} + \delta X_{i,2} = 0$$
at first order
$$\Delta\Phi^{(1)} = \sum_{i=1}^{r}\left[\left(\pdv{\Phi}{X_i}\right)_1 -
\left(\pdv[]{\Phi}{X_i}\right)_2\right]\delta X_{i,1}\geq0$$
As before, the only way to ensure this condition is
$$\left(\pdv[]{\Phi}{X_i}\right)_1 = \left(\pdv[]{\Phi}{X_i}\right)_2 
\leftrightarrow I_{i,1} = I_{i,2}$$
with finer partitions this means that the intensive variables conjugated to the
extensive control variables are homogeneous in the system, at leas at first order.

To get a non-trivial result, consider the second order of fluctuation of a single
variable $X_i$ (i.e. considering the fluctuations independent):
$$\Delta\Phi^{(2)} = \frac{1}{2}\left[\left(\pdv[2]{\Phi}{X_i}\right)_1 +
\left(\pdv[2]{\Phi}{X_i}\right)_2\right](\delta X_{i,1})^2\geq0	$$
where the minus sign was destroyed by the square; this statement is true only if
$$\pdv[2]{\Phi}{X_i} = \pdv[]{I_i}{X_i}\geq0$$
from this follows that
$$\pdv[]{X_i}{I_i}\geq0$$
and, remembering the Legendre transform
$$\Phi(...,X_i,...)\rightarrow\tilde\Phi(...,I_i...) \leftrightarrow X_i =
-\pdv[]{\tilde\Phi}{I_i}$$
this gives
$$\pdv[2]{\tilde \Phi}{I_i}\leq0$$

This gives a geometric condition on the TD potentials at equilibrium: they are either
a concave function of intensive variables or convex functions of extensive
variables (or vice-versa \textbf{CHECK!!!!!!!})
\textbf{Example} Consider the Helmholtz free energy: the stability condition
implies that
$$\pdv[2]{F}{T}\leq0$$
but
$$C_V = -T\pdv[2]{F}{T}\rightarrow C_V\geq0$$
and also
$$\underbrace{\pdv[2]{F}{V}}_{\geq0}=-\pdv[]{p}{V} = -\frac{1}{\pdv[]{V}{p}} =
\frac{1}{\chi_T}\geq0 \leftrightarrow \chi_T\geq0$$
\textbf{Example} Consider the Gibbs energy $G(T,p,N)$:
$$\pdv[2]{G}{T}\leq0$$
and the $C_p$ definition
$$C_p = -T\pdv[2]{G}{T}\geq0.$$
Again the $\chi_T$ definition is
$$\chi_T = -\underbrace{\pdv[2]{G}{p}}_{\leq0}\geq0$$
which is consistent with the previous example.

From the relation
$$C_p = C_V + \frac{T\alpha_p^2}{\chi_T}$$
and the stability conditions discussed it follows that, in general
$$C_p\geq C_V$$
\subsection{Multi-component systems}
Some TD systems have different components mixed (e.g. the atmosphere), indexed with
$a=1,2...c$ so the chemical potential is not unique: $\mu_a,N_a$ and these potentials
depend separately from those variables (which vars???)
$$\dd E = T\dd S -p\dd V + \sum_{a=1}^{c}\mu_a \dd N_a + y\dd X$$
Each potential has its own equilibrium condition $\mu_{a,1} = \mu_{a,2}$ and
intensive variables depend only on $c-1$ ratios
$$\nu_a = \frac{N_a}{N}$$
with the constraint
$$\sum_{a=1}^{c}N_a =N\leftrightarrow\sum_{a=1}^{c}\nu_a=1$$

\subsection{Phase transitions}
This topic is very difficult to treat mathematically and one of the most important
results of statistical mechanics is the description of this phenomenon.

What are the conditions of coexistence of different phases of matter in equilibrium
(e.g. ice and water)?\\
At equilibrium, the state of the TD system is characterized by $E,V,N$: it is not
necessary that at equilibrium state these variables are homogeneous.

Each different phase (each defined by its state equation) of the system has its set
of variables $p_i,T_i,\mu_i$.

If the state is equilibrium, the intensive variables must be homogeneous:
$$p_1=p_2,\quad T_1=T_2,\quad\mu_1=\mu_2$$
it is then natural to take $p,T$ as control variables, making $\mu$ a dependent
variable; since the phases have different equations of state, the potentials will
be different: for the phase $\alpha$
$$G_\alpha(T,p,N_\alpha) = \mu_\alpha N_\alpha \rightarrow \mu_\alpha =
\frac{G_\alpha}{N_\alpha}.$$
The phases can only coexist on \textit{curves of coexistence} in the state space:
for generic values of $p,T$, $\mu_\alpha\neq\mu_\beta$ and so the phases do not
coexist and, considering $N$ constant, the system "selects" the phase which minimizes
$$\mu=\frac{G}{N}.$$
So phases $\alpha$ and $\beta$ coexist if and only if the values of $p,T$ are so that
$$\mu_\alpha(p,T) = \mu_\beta(p,T)$$
which typically is a curve $p(T)$ in the state space. When this curve is crossed, a
phase transition occurs and it can be of two natures.
\subsubsection{First order transitions}   
The coexistence curve is determined by the intersection of the surfaces
$\mu_\alpha(T,p)$ and $\mu_\beta(T,p)$. In general, the value of the derivatives
differ at the intersection point(s): fixing $p$ they are
$$\left(\pdv[]{\mu_\alpha}{T}\right)_{p} = \frac{1}{N_\alpha}\left(\pdv[]{G}{T}
\right)_p = -\frac{S_\alpha}{N_\alpha}:=-s_\alpha$$
which is the entropy per particle, so in general
$$s_\alpha\neq s_\beta$$
Suppose $s_\alpha >s_\beta$: for $T=T(p)-\varepsilon$, $S_{(-)} = Ns_\beta$,
while for $T=T(p)+\varepsilon$, $S_{(+)} = Ns_\alpha$; in this case the
transition from $\beta$ to $\alpha$ makes the entropy jump of
$$\Delta S = N(s_\alpha-s_\beta)>0$$
this is nothing but the effect of latent heat, the heat to be furnished to the
system just to do a phase transition
$$L=T(p)\Delta S.$$
From an enthalpy p.o.v,
$$H(S,p,N) = G+TS$$
since $\mu_\alpha = \mu_\beta$, $\Delta G = 0$ and so
$$\Delta H = \Delta(TS) = T\Delta S = L.$$\\

This is similar for fixed temperature transitions, with derivative
$$\left(\pdv[]{\mu}{p}\right)_{T,N} = \frac{V}{N}=v$$
which is the volume per particle; so in general $v_\alpha\neq v_\beta$ and the
volume of the system changes during the transition
\textit{even at constant pressure} (e.g. ice melting).
